% ============ 04. Matriz D ============
\section{Matriz detallada D1--D5: qué hay, qué falta, qué escribir}

\subsection{D1 · Problema, contexto y estado actual (12\,\%, bloqueante)}

\textbf{Nivel: En desarrollo} (riesgo de Insuficiente si se entrega sin
redactar). Lo que \emph{sí} existe: abundantes cifras propias verificadas
(77\,237 registros declarados / 30\,086 investigadores únicos; 51\,\%
Bogotá+Antioquia; 24\,\% mujeres en Ingeniería vs.\ 48\,\% en ciencias médicas;
afro 3$\times$, indígena 7,9$\times$, discapacidad 8$\times$ vs.\ DANE; 22
atípicos de edad; 10/16 scripts reproducibles). Lo que \emph{no} existe aún:
enunciado del problema con \textbf{actor, alcance, frecuencia y consecuencia};
la \textbf{pregunta estadística respondible}; y (modalidad datos abiertos) la
\textbf{fuente externa} que documente la magnitud del problema en Colombia y a
quién le importa.

\begin{tcolorbox}[nota]
\textbf{Para llegar a Ejemplar hay que escribir (1\,½--2 págs.):} quién sufre el
problema (sistema de reconocimiento / investigadores / política de CTI), desde
cuándo y con qué frecuencia (convocatorias 2013--2021), cómo se resuelve hoy
sin intervención (con una cifra propia), qué pasa si nadie hace nada, y la
traducción a una pregunta estadística respondible (p.~ej., ``¿el padrón de
investigadores reconocidos de MinCiencias (2013--2021) es fiable, comparable
entre convocatorias y equitativo?''), apoyada en una fuente que documente la
magnitud del problema en Colombia.
\end{tcolorbox}

\subsection{D2 · Estado del arte y rigor bibliográfico (9\,\%)}

\textbf{Nivel: En desarrollo} (el activo más maduro, con fallas donde la rúbrica
mira primero). El \texttt{estado\_del\_arte.pdf} (20 págs., 39 referencias)
está organizado por enfoque, distingue mundo/Colombia y fue validado con
\textbf{0 referencias fabricadas} (confirmado por dos agentes). Fallas frente a
la rúbrica: (i) la citación es \textbf{autor-año (plainnat), NO APA 7}
(corchetes, ``and'', sin DOI); (ii) no contrasta explícitamente \textbf{dos
aproximaciones señalando en qué difieren sus supuestos}; (iii) no declara el
\textbf{vacío} que atiende el proyecto nuevo; (iv) referencias cruzadas internas
del PDF corregidas en esta revisión (sección 13.2).

\begin{tcolorbox}[nota]
\textbf{Para Ejemplar:} sintetizar el SotA en 2--3 págs. del anteproyecto
(6--7 enfoques), añadir un contraste explícito (p.~ej., modelos de Markov vs.\
supervivencia para retención; HHI vs.\ Gini/Theil), declarar el vacío (falta de
una auditoría de calidad y equidad reproducible del registro colombiano con
cruce de productividad), y \textbf{convertir la bibliografía a APA 7} (~18--25
referencias citadas, con DOI verificados).
\end{tcolorbox}

\subsection{D3 · Objetivos (8\,\%, bloqueante)}

\textbf{Nivel: En desarrollo} (riesgo de Insuficiente si se copian los OE
actuales). El informe base (\S3.1--3.2) documenta la consultoría \emph{ya
ejecutada} con verbos de actividad (\emph{ampliar}, \emph{robustecer},
\emph{reparar}), sin evidencia de cumplimiento medible y sin un problema D1 al
que responder. Los objetivos no son heredables tal cual: \textbf{deben
reescribirse como logros del anteproyecto}, alineados al problema elegido.

\begin{tcolorbox}[nota]
\textbf{Para Ejemplar:} un objetivo general que responda exactamente al
problema de D1 (ni más ancho ni más angosto) y \textbf{entre 3 y 5 específicos}
que lo agoten sin solaparse, cada uno con verbo verificable (p.~ej.,
\emph{cuantificar}, \emph{estimar}, \emph{validar}, \emph{entregar}) y la
evidencia con que se declarará cumplido. Los OE1--OE3 del informe son una base
conceptual que debe reformularse, no copiarse.
\end{tcolorbox}

\subsection{D4 · Datos y aproximación metodológica preliminar (6\,\%)}

\textbf{Nivel: Competente} (único criterio por encima de ``En desarrollo'').
Procedencia, periodo, unidad de observación, volumen, acceso y diccionario de
variables están verificados en \texttt{datos/catalogo.yaml} y el README
(Socrata datos.gov.co: \texttt{bqtm-4y2h} y \texttt{33dq-ab5a}). Para
Ejemplar faltan cuatro piezas: \textbf{licencia declarada}, \textbf{plan B si
el dato no llega}, una \textbf{tabla objetivo $\rightarrow$ variables
$\rightarrow$ método}, y el argumento de \textbf{``no está ya resuelto y
publicado con esos mismos datos''}.

\subsection{D5 · Alcance, viabilidad y plan de trabajo (5\,\%)}

\textbf{Nivel: En desarrollo}. El informe base (\S4) tiene exclusiones
declaradas con razón y 12 riesgos con mitigaciones — base sólida. Falta el
\textbf{cronograma por hitos con entregables verificables y holgura} (12
semanas, 2 de margen) y el \textbf{presupuesto de horas} real por integrante.
